\documentclass[tikz,border=5]{standalone}
\usepackage{amsmath}
\usetikzlibrary{calc}
\begin{document}       
\tikzset{
    relative to node/.style={
        shift={(#1.center)},
        x={(#1.east)},
        y={(#1.north)},
    }
}
\begin{tikzpicture}
    \node[above right,inner sep = 0] (image) at (0,0) {
        \includegraphics[width=\linewidth]{../figs/example-fig-beta-vT.pdf}
    };
    \begin{scope}[relative to node=image]
        % \foreach \ratio in {-1,-0.9,...,1}{
        %    \draw[help lines,purple] (-1,\ratio) -- (1,\ratio);
        %    \draw[help lines,purple] (\ratio,-1) -- (\ratio,1);
        % }
        % \draw (-1,0) -- (1,0);
        % \draw (0,-1) -- (0,1);
        
        \node[anchor=center,blue] at (0.25,-0.47) {$\Delta T\approx 2.0^\circ\text{C}$};

        \node[anchor=center,blue,rotate=90] at (0.5,0.0) {$\Delta v\approx2.1\,\ell/\text{kg}$};
      
        \coordinate (pvT) at (0.26,0.03);
        \draw[latex-, thick,shorten <=3pt,color=blue] (pvT) to[bend right=25] ++(-0.2,0.2) node[anchor=east] {$\left.\begin{aligned}
            p&=225\,\text{bar}\\
            T&=376^\circ\text{C}\\
            v&\approx 4.0\,\ell/\text{kg}
        \end{aligned}\right\}$};
        
        
    \end{scope}
\end{tikzpicture}
\end{document}
